\scanpage{14}
\uline{\unsure{Résolutions|Résolution} standard et théories relatives.}

\diagram{\begin{tikzcd}[ampersand replacement=\&]
 C'\arrow[r,shift left,"G"] \& C\arrow[l,shift left,"F"]
\end{tikzcd}}
% slanted
\margin{\uline{\unsure{Dont|Droit|Donné}}}
% slanted
\margin{\unsure{Godement|G-préfaisc.|G-faisc.}}
\[
 \mathrm{Hom}(F(X),X')\xrightarrow{\sim}\mathrm{Hom}(X,G(X'))
 \qquad F\quad G
\]
Posons
\[
 E=GF:C\longrightarrow C
\]
d'où
\[
 \varphi:\mathrm{id}_C=E^0\longrightarrow E,\qquad \chi:E^2\longrightarrow E
\]
satisfaisant les relations de Godement, d'où un objet \inserted{\textcircled{\scriptsize co}}\uline{simplicial} \deleted{dans} $\widehat E$ dans \uline{Hom} $(C,C)$ \deleted{\unsure{$\varphi$ et $\chi$ où}}.
\[
 \widehat E^{\,n}=E^{n+1}\ \deleted{\unsure{n}}\ \ldots\ldots\quad
 \text{et \uline{augmentation}}\quad X\longrightarrow\widehat E(X)
\]
(On devrait trouver une interprétation naturelle par théorie de la descente). On pose de \textquotedbl
\[
 E'=FG:C\longrightarrow C
\]
d'où
\[
 \varphi':E'\longrightarrow\mathrm{id}_{C'}=E'^{\,0},\qquad
 \chi':E'\mathrel{\unsure{-|\to}}E'^{\,2}
\]
satisfaisant aux relations de Godement, d'où un objet \uline{simplicial} $\widehat E'$ de \uline{Hom} $(C',C')$ :
\[
 \widehat E'^{\deleted{\unsure{n}}}_{n}=E'^{\,n+1}\qquad\text{---}
\]
Notons les \deleted{\unsure{isomorphismes} \unsure{fonctoriels|simpliciaux}} \inserted{\unsure{résultats suivants}}
\[
 \begin{aligned}
 \text{(i)}\quad &G(X')\mathrel{\substack{\sim\\[-.7ex]\sim\\[-.7ex]\sim}}\widehat E(G(X'))\ \deleted{\unsure{\simeq0}}
       &&\text{pour tt }X'\in\mathrm{ob}\,C'\\
 \text{(ii)}\quad &F(X)\approxeq F(\widehat E(X))\ \deleted{\unsure{\simeq0}}
       &&\text{pour tt }X\in\mathrm{ob}\,C
 \end{aligned}
\]
La deuxième relation est vraie plus gén. si on a un foncteur $T:C\longrightarrow D$ et un hom \unsure{fonctoriel}
\[
 h:T\circ E\longrightarrow T
\]
tel que
\[
 h\circ(T*\varphi)=\mathrm{id}_T
\]
alors
\[
 T\widehat E(X)\approxeq 0\qquad\text{pour tt }X\in\mathrm{ob}\,C
\]
[cf Godement]. La formule (i) se vérifie sans doute de \textquotedbl par des formules explicites.

\uline{Application} Supposons $C,C'$ abéliennes, $F,\unsure{G|C}$ additifs, \deleted{\illegible{2}} \deleted{\illegible{1}}. a) Si $F$ est \uline{exact} et \uline{fidèle}, alors $\widehat E(X)$ est une \uline{résolution} de $X$ pour tt $X$.

b) Si $\Gamma$ est un foncteur
% archive sheet number: 13
